\section{Discussion}
\label{sec:discussion}

Every analysed detector passes the known-groups check \citep{hattie1984procedures}, and on Qwen so
does every random direction tuned like the probe. The check therefore cannot tell these
detectors apart, whereas agreement within cue conditions (post hoc) can. Self-report and
third-person reach most of $\kappa_{\max}$, which suggests that they respond to largely the same
evidence at different rates. The Qwen probe, by contrast, agrees with neither.

Three practices follow. First, report each detector's hit and false-alarm rates beside $d'$ and
beside any threshold rule. Second, report $\kappa$ beside $\kappa_{\max}$, both computed within
each cue condition and any other design variable that shifts positive rates. Third, compare a
fitted detector with a null from its own fitting procedure \citep{hewitt2019designing}, as
Figure~\ref{fig:probe-null} does.

Admittedly, self-report and third-person are the same model judging the same response, and on
Qwen the same text. Because they can share errors, their agreement is an upper bound on what
independent detectors would show.
